\section{Software: Arcane Framework}
\label{sec:Arcane Framework:software}



\begin{FicheMetadataTable}[Arcane Framework Information]
\textbf{Consortium} & \begin{tabular}[t]{@{}p{\linewidth}@{}}
CEA\\
IFPEN\\
\end{tabular} \\
        \textbf{Exa-MA Partners} & \begin{tabular}[t]{@{}p{\linewidth}@{}}
CEA\\
\end{tabular} \\
        \textbf{Contact Emails} & \begin{tabular}[t]{@{}p{\linewidth}@{}}
lydie.grospellier@cea.fr\\
\end{tabular} \\
        \textbf{Supported Architectures} & \begin{tabular}[t]{@{}p{\linewidth}@{}}
CPU Only\\
CPU or GPU\\
\end{tabular} \\
        \textbf{Repository} & \href{https://github.com/arcaneframework/framework}{https://github.com/arcaneframework/framework} \\
        \textbf{License} & \begin{tabular}[t]{@{}p{\linewidth}@{}}
OSS: apache-2\\
\end{tabular} \\
        \textbf{Bottlenecks roadmap} & \begin{tabular}[t]{@{}p{\linewidth}@{}}
B10 - Scientific Productivity\\
B11 - Reproducibility and Replicability of Computation\\
B6 - Data Management\\
B7 - Exascale Algorithms\\
\end{tabular} \\
\end{FicheMetadataTable}

\subsection{Software summary}
\label{sec:Arcane Framework:summary}
Detailed overview not available.



\subsection{Purpose}
\label{sec:Arcane Framework:purpose}
Purpose not available.

\subsection{Programming and Computational Environment}
\label{sec::Arcane Framework:environment_capabilities}


The following table summarizes these aspects for Arcane Framework, providing a  view of its programming and computational capabilities.

\begin{FicheEnvironmentTable}[Arcane Framework programming and computational environment]
Languages  & \begin{tabular}[t]{@{}p{\linewidth}@{}}
C++\\
C\#\\
\end{tabular} & Programming languages and language standards supported by the software \\
        Parallelism  & \begin{tabular}[t]{@{}p{\linewidth}@{}}
GPU\\
MPI\\
Multithread\\
\end{tabular} & Parallel computing methods and frameworks utilized by the software.\\
        Data Formats  & \begin{tabular}[t]{@{}p{\linewidth}@{}}
Ensight\\
HDF5\\
JSON\\
XML\\
\end{tabular} & Data formats that the software can handle or produce.\\
        Resilience  & \begin{tabular}[t]{@{}p{\linewidth}@{}}
Checkpoint restart\\
\end{tabular} & Fault tolerance and recovery mechanisms employed by the software.\\
        DevOps & \begin{tabular}[t]{@{}p{\linewidth}@{}}
Continuous Integration\\
\end{tabular} & Outlines the development and operational practices including continuous integration, containerization, and testing methodologies.  \\
        Packaging  & \begin{tabular}[t]{@{}p{\linewidth}@{}}
None\\
\end{tabular} & Software packaging and distribution.\\
        Testing  & \begin{tabular}[t]{@{}p{\linewidth}@{}}
Unit\\
Verification\\
\end{tabular} & Testing methodologies employed to ensure software quality and correctness.\\
        Containerization  & \begin{tabular}[t]{@{}p{\linewidth}@{}}
None\\
\end{tabular} & Container technologies used to package and deploy the software.\\
        Interfaces  & \begin{tabular}[t]{@{}p{\linewidth}@{}}
None\\
\end{tabular} & List of software Arcane Framework has interfaces with.\\
\end{FicheEnvironmentTable}



\subsection{Mathematics}
\label{sec:Arcane Framework:mathematics}
Mathematics not available.

In this section, provide a summary the mathematics used in the software.


\subsection{Relevant Publications}
\label{sec:Arcane Framework:publications}

Here is a list of relevant publications related to the software:


\subsection{Acknowledgements}
\label{sec::Arcane Framework:acknowledgements}

The software has been developed with the support of the following funding agencies and institutions: 




Acknowledgements not available.


